\documentclass[11pt,a4paper]{article}
\usepackage[utf8]{inputenc}
\usepackage[T1]{fontenc}
\usepackage[brazil]{babel}
\usepackage[a4paper,margin=2.3cm]{geometry}
\usepackage{enumitem}
\usepackage{hyperref}
\usepackage{xcolor}
\setlength{\parindent}{0pt}
\setlength{\parskip}{5pt}
\sloppy
\setlength{\emergencystretch}{3em}
\hypersetup{colorlinks=true,urlcolor=blue,linkcolor=black}
\newcommand{\cabecalho}[2]{\begin{center}{\small \textbf{Sprint 2 -- #1}}\\[1mm]\textit{Inventário de Ativos e Vulnerabilidades}\\Matheus Moreli Durão Silva -- 12621CBS217\\Cibersegurança -- Bacharelado -- Noturno\end{center}\hrule\vspace{4mm}}
\newcommand{\bloco}[1]{\par\vspace{2mm}\textbf{#1}\par}
\newcommand{\codigo}[1]{\texttt{#1}}

\begin{document}\cabecalho{S2\_06 -- Critérios de Aceitação do CRUD}{}
\bloco{Cenários observáveis} CA01/RF01: dado um ID livre e campos válidos, quando cadastrar, então o ativo aparece na consulta por ID. CA02/RF01: dado um ID já usado, quando cadastrar, então a base não muda e a mensagem informa duplicidade. CA03/RF04: dado o ativo 10, quando atualizar responsável, então o nome novo aparece e o ID continua 10. CA04/RF05: dado o ativo 10, quando a confirmação for diferente de SIM, então ele permanece na base. CA05/RF06: dado um ativo existente e CVE/CVSS válidos, quando cadastrar a vulnerabilidade, então ela aparece na consulta pelo CVE.
\bloco{Cenário de erro} CA06/RN02: dado CVSS 11, quando cadastrar vulnerabilidade, então o cadastro é recusado sem alterar a lista. CA07/RF02: dado ID inexistente, quando consultar, então o resultado é ``não encontrado'', sem erro técnico.
\bloco{Como viram testes} Os cenários viraram métodos que começam com \codigo{testar\_} em \codigo{test\_sprint2.py}. O esperado é conferido por \codigo{assertEqual}, \codigo{assertIn} ou \codigo{assertRaises}.
\bloco{Checklist atendido} Contexto, ação e resultado; sucesso e falha; resultado visível; requisito indicado; possibilidade de automação.
\bloco{Matriz completa de cenários} O arquivo versionado \codigo{docs/criterios\_aceitacao.md} registra CA01--CA16 com requisito, Dado, Quando e Então. O CRUD tem ao menos dois casos por operação: cadastro válido/duplicado; consulta existente/inexistente; atualização válida/nome já usado; exclusão cancelada/confirmada com limpeza de dependências. Os resultados são observáveis no dicionário de objetos e no JSON recarregado. A revisão por outro integrante foi solicitada pelo card e ainda depende de registro humano; não é declarada como concluída.
\end{document}
